% =====================================================================
\section{Limitações, ameaças à validade e dificuldades encontradas}
\label{sec:limitacoes}

\subsection{Limitações do que foi provado}

As garantias obtidas neste trabalho são condicionais, e a lista abaixo delimita exatamente a
que condições:

\begin{itemize}[leftmargin=1.4cm]
  \item \textbf{Horizonte finito.} Todos os vereditos de BMC valem até o limite de
  desenrolamento indicado. A verificação em malha fechada cobre dois passos; o
  \textit{harness} de cinquenta passos não foi concluído e a medição mostra que o custo
  torna-o inviável nesta máquina.
  \item \textbf{Indecisão sob faixas estreitas.} Abaixo de 1,25\,rad/s, a propriedade de dois
  passos permanece indecisa em 300\,s. Não há, portanto, faixa de operação em que o
  controlador esteja provado seguro.
  \item \textbf{Nenhum dos dois métodos decide o controlador real.} O PDR não convergiu em
  1802\,s e o BMC alcançou cinco \textit{frames} em 901\,s.
  \item \textbf{Sobreaproximação nas ativações.} Os quatro vereditos de blocos FFN valem para
  a relaxação intervalar injetada. Transferi-los para GeLU e SiLU reais exige que esses
  limites sejam conservadores, o que não foi provado --- o \textit{encoding} não representa a
  relação funcional exata.
  \item \textbf{Aproximação linear da sigmoide.} Onde a forma funcional foi mantida, a
  substituição por $\operatorname{clip}(0{,}25z + 0{,}5, 0, 1)$ altera a função implantada; o
  veredito vale para a aproximação.
  \item \textbf{Perfis de caos com faixa contínua não decidem} sob \texttt{--floatbv}; apenas
  o perfil discreto de impulso decide, em 27\,s.
  \item \textbf{Verificação a partir de Python não realizada.} O binário disponível não possui
  o \textit{frontend} Python compilado.
  \item \textbf{Biblioteca de terceiros não verificada.} O bloqueio é a ausência de
  \texttt{--std} na versão 6.8.0, necessária para o construtor delegante de C++11.
\end{itemize}

\subsection{Ameaças à validade}

\subsubsection{Validade de construção: a propriedade vácua}

A ameaça mais séria identificada no ciclo não é de desempenho, mas de significado. Uma
propriedade que decide sem envolver o artefato produz uma prova verdadeira sobre um objeto
que não interessa, e não emite nenhum sinal de erro. Duas das três propriedades da malha
fechada estavam nessa condição, e foram descobertas apenas pelo teste explícito de
dependência descrito na Seção~\ref{sub:vacuidade}. Não há razão para supor que o repositório
esteja hoje livre de casos análogos; o que existe é um procedimento para detectá-los, e ele
deve ser aplicado a toda propriedade nova.

\subsubsection{Validade interna: fidelidade das traduções}

Três traduções sustentam os resultados e nenhuma delas está provada correta. A quantização
Q8.8 do ator coincide bit a bit entre três implementações, o que é evidência forte de
consistência interna, mas não de fidelidade ao modelo treinado em ponto flutuante --- a
caracterização de erro mostra pior caso de 7,76\,N. A tradução para Verilog foi validada por
teste diferencial em doze estados: isso detecta regressão, não equivalência. E os
\textit{stubs} de blocos de \textit{transformer} representam fatias reduzidas, não os módulos
de produção.

\subsubsection{Validade externa: escala e generalização}

Os experimentos foram conduzidos em uma única máquina de quatro núcleos, com um binário
específico do verificador. As barreiras medidas --- $N{=}4$ no GEMM, $K{=}4$ na malha fechada,
cinco \textit{frames} no ABC --- são propriedades daquele ambiente e daqueles
\textit{harnesses}, e não constantes do método. Extrapolá-las seria repetir o defeito que a
auditoria do projeto encontrou e corrigiu.

\subsubsection{Validade de conclusão: procedência dos números}

Este relatório reporta apenas números com log bruto versionado, com uma exceção declarada: a
execução histórica de PDR sobre o ator real, cuja saída bruta se perdeu antes da correção do
\textit{script}. Ela é apresentada como medição histórica do invólucro, e não como resultado
plenamente reproduzível.

\subsection{Dificuldades encontradas ao longo do ciclo}

Registram-se as dificuldades que mais consumiram tempo, na expectativa de que sejam úteis a
quem retomar o trabalho:

\begin{enumerate}[leftmargin=1.4cm]
  \item \textbf{Defeitos silenciosos na instrumentação.} Vários experimentos produziam números
  plausíveis sem medir nada: um parâmetro de dimensão que não chegava ao código compilado, uma
  \textit{flag} que fazia o verificador emitir a fórmula em vez de resolvê-la, um portão de
  verificação comentado em um orquestrador. Nenhum deles emitia erro. Foram encontrados apenas
  por auditoria dirigida, e não por execução da suíte.
  \item \textbf{Suíte de testes verde sem verificar.} A suíte chegou a passar sem exercitar o
  verificador, porque o resultado era reduzido a um booleano que também é falso quando nada
  executa. A reescrita do invólucro e a suíte de regressão de casos-limite foram respostas
  diretas a isso.
  \item \textbf{Divergência entre binário e árvore de código.} O binário versionado é a versão
  6.8.0 e a árvore presente no repositório é a 8.0.0, não compilada. \textit{Flags} da 8.0.0
  são rejeitadas pela 6.8.0 com código 64, facilmente confundível com propriedade violada.
  Compilar a 8.0.0 leva de uma a quatro horas, o que tornou a alternância custosa.
  \item \textbf{Escolha de limites de tempo.} O episódio do PID mostra o custo de um limite
  arbitrário: 200\,s produziram \textit{timeout} para uma propriedade que é verdadeira e se
  prova em pouco mais que isso.
  \item \textbf{Submódulo Git inconsistente.} Um \textit{gitlink} sem entrada correspondente em
  \texttt{.gitmodules} fazia o \textit{checkout} automatizado terminar com erro. O conteúdo foi
  recuperado --- e não descartado --- após identificar o repositório e o \textit{commit}
  corretos.
  \item \textbf{Peso do repositório.} Material de terceiros vendorizado responde por
  308\,MB dos 322\,MB do diretório do projeto, o que penaliza clonagem e integração contínua.
  A decisão sobre o que preservar permanece em aberto.
\end{enumerate}
